\documentclass[11pt]{article}
\usepackage[left=1.25in,top=1.0in,twoside=false,bindingoffset=0in]{geometry}

\usepackage{fontspec}
\setmainfont{Latin Modern Sans}
%\setmainfont{TeX Gyre Termes}

\begin{document}

\begin{center}
\LARGE Chouette Rules
\end{center}

\begin{center}
  (Under construction: \today)
\end{center}

\subsection*{Starting}

\begin{enumerate}

\item Write the agreed upon per-point stake on the score sheet.

\item By default, the Jacoby rule is in effect, and no Beavers, Partners,
  or Extras are allowed.  To change any of these, indicate your
  choices on the score sheet.  Sections below summarize the specific
  rules that govern Partners and Extras.

\item A designated person or player maintains the score sheet, an example of
  which is attached, and maintains the order of the rotation.

\item The players roll dice to determine the initial rotation.
The first player in the rotation becomes the initial Box.
The second player in the rotation becomes the initial Captain.

\item All the remaining players comprise the initial Crew and play on the
  Captain's side against the Box.  There is a single game at the board,
  with the Box and the Captain moving the checkers, but each of the Crew
  and the Captain has his own independent doubling cube against the Box.
  Conversely, the Box uses those independent doubling cubes against each
  of the Crew and the Captain, and the games are scored as if they were
  independent games with the doubling cube multiplying the above-agreed
  stake, with gammons and backgammons counting according to Jacoby if in
  effect.  All participants -- Box, Captain, and each of the Crew -- make
  independent cube decisions.

\item The score sheet must be updated and verified before each game begins.


\end{enumerate}

\subsection*{Game Mechanics}

\begin{enumerate}
\item The Captain moves the checkers on his behalf and on behalf of
  the Crew, and has final say over any differing opinions about the
  best checker play.

\item A crew member may Consult with the Captain only if both his cube
  and the Captain's are in the game (not centered).  In particular,
  before the Captain's cube is in the game, the Crew may not consult
  with the Captain (or each other), and their role is to observe the
  game and to make their own independent cube decisions.  Nobody whose
  cube is not in the game can consult in any manner.

\item Each of the cubes operates as in a normal heads-up game of
  backgammon, with one exception: the Box must double the Captain on
  any initial double if the Captain's cube is still centered.  In
  other words, he may not double a Crew member with a centered cube
  without also doubling the Captain if the Captain's cube is still
  centered.

\item Subject to the above constraint, the Box is free to
  independently double, redouble, take, or pass each cube
  independently of all of the others.  He may double some but not
  others, he may take some cubes and drop some, etc.
\end{enumerate}

\subsection*{Rotation}

\begin{enumerate}
\item If the Captain loses, he goes to the bottom of the rotation and the
next player in the rotation who is still in the game takes his seat
as Acting Captain until the end of the game, or until he is cubed out,
in which case the next person in the rotation who is still in the
game becomes Acting Captain, etc.

\item If the Box loses to the Captain or fails to at least break even
against his opponents as a whole that game, he loses the Box and goes
to the bottom of the rotation at the completion of the game.

\item The order of the rotation does not change otherwise.  In particular,
winning or losing as Acting Captain does not change one's place in the
rotation.

\item It is the order of the rotation at the start of a game that always
  determines who is Box (first in rotation) and who is Captain (second in
  rotation) and who is Crew (everybody else).

As a consequence of the above rules:
\begin{itemize}
  \item If the Captain beats the Box, he will always be the Box next game.

  \item If the Box wins against the Captain and at least breaks even, he
    will remain Box next game.

  \item If the Captain loses and the Box fails to break even, both go to
    the bottom of the rotation (with the former Box last).
\end{itemize}

\end{enumerate}

\subsection*{Consulting}

\begin{enumerate}
\item When Consultation between a Crew member and the Captain is
  permitted by the above rules, such discussion may be requested by
  the Captain or offered by the Crew member.

\item
  These discussions must be about checker play, not cube decisions, so
  no such discussion can occur at all until after the Captain has rolled
  the dice and there are checkers to play.

\item When discussing available checker plays, it is permissible to
  discuss the merits of each, including how future possible cube
  action affects those plays.

\end{enumerate}

\subsection*{Cube Order}

\begin{enumerate}
\item Even when consulting is not allowed, a member of the Crew may
  ask the Captain to hold his roll so he can consider whether or not
  he will double.

\item Each of the Crew, and the Captain, is permitted to announce his
  decision to double, or to announce his decision not to double, or to
  take or drop an offered cube, in any order.  However, each member of
  the Crew has a right to hear the final cube decision of every player
  on his side higher in the rotation than he is, before he makes his
  own final cube decision.

\item Neither Captain nor any Crew member may announce or indicate the
  reason for any cube decision he makes, only the decision itself.
\end{enumerate}


\subsection*{Cube Proxy}

\begin{enumerate}
\item There is a limited exception to the rule that the Box's
  opponents make independent cube decisions: if a Crew member needs to
  leave the chouette for a short time, he may ask another Crew member,
  or the Captain, to handle his cube while he is away.  The Captain
  handles his cube if no other proxy was obtained.  The proxy expires
  at the end of the game, or when the player returns.  If the player
  has not returned by then, he is treated as having temporarily left
  the chouette, as described in the next section.

\item There is no good way to handle the Captain or the Box leaving in
  the middle of a game, except to wait, which obviously should be
  avoided if possible.

\end{enumerate}

\subsection*{Entering \& Leaving}

\begin{enumerate}
\item A new player can join the chouette by agreement of the existing
  players.  He enters the rotation at the bottom, but only after any
  existing game is completed and all the existing players have taken
  their new positions in the rotation and their new roles in the
  upcoming game.  Then the new player is added at the bottom as a
  member of the Crew.

\item A player may temporarily leave the chouette, for
example to play a tournament match.  In such a case, the leaving
player is taken out of the rotation and put to the side.  When he
re-enters, he enters the same as a new player, at the end of the
rotation, but with his score maintained from when he left.

\item Any player may permanently leave a chouette after the completion of
any game. He either cashes out his winnings from other players or pays
other players to cash out, depending on whether his score is positive
or negative.  All players are expected to fully settle at the end of
the chouette, or when they leave the chouette.
\end{enumerate}

\subsection*{Partners}

This section describes the rules that apply only when the Partners option was
written on the score sheet at the beginning of the chouette.

\begin{enumerate}
\item Partners are generally used only in a chouette of six or more
  players.

\item Before a game, the Box may invite the player next in the rotation
  after the Captain to be his Partner.  He may decline without losing
  his place, in which case the next player in the rotation is invited,
  and so on.  If everyone declines, the Box plays alone.

\item The Partner plays on the Box's side, against the Captain and the
  Crew, for that game only.  He keeps his place in the rotation as
  though he were still Crew, so the rotation rules above apply
  unchanged: the order of the rotation at the start of each game
  decides who is Box and who is Captain, and the Box may invite a
  Partner again.

\item The Box moves the checkers and makes the decisions for the pair;
  the Partner does not handle the dice or checkers.  The Box and
  Partner may discuss cube decisions at any time, but not specific
  rolls or checker plays.  They may discuss checker play freely only
  after two cubes have been turned, whether taken or passed.

\item The pair share their points equally: their total result against all
  of their opponents is divided between them, with an odd point going to
  the Box when they win and being lost by the Box when they lose.  The Box
  keeps the Box on the same terms as a lone Box, judged on the pair's
  combined result.

\item Private settlements between Box and Partner change only their own
  shares.  They do not affect the rotation or who keeps the Box, except
  that a player left with no financial interest in the rest of the game
  may no longer be consulted.
\end{enumerate}

\subsection*{Extras}

\noindent TBD.

\end{document}
